\documentclass[11pt]{amsart}
\usepackage{amsmath,amssymb,mathtools}
\usepackage[hidelinks]{hyperref}
\hypersetup{pdftitle={Exact Local Dimension of the Brent Variety V(4,4,4 | 49) at the Recursive Strassen-Square Point via Quadratic Obstructions},pdfauthor={Edson Soares}}
\addtolength{\textwidth}{18pt}
\addtolength{\oddsidemargin}{-9pt}
\addtolength{\evensidemargin}{-9pt}
\addtolength{\textheight}{20pt}
\raggedbottom
\newtheorem{theorem}{Theorem}
\newtheorem{lemma}[theorem]{Lemma}
\newcommand{\F}{\mathbb F}
\newcommand{\Q}{\mathbb Q}
\newcommand{\Z}{\mathbb Z}
\newcommand{\Br}{\mathcal B}
\newcommand{\rk}{\operatorname{rank}}
\newcommand{\im}{\operatorname{im}}
\title[Local dimension at the Strassen square]{Exact Local Dimension of the Brent Variety
\texorpdfstring{$V(4,4,4\mid49)$}{V(4,4,4 | 49)} at the Recursive
Strassen-Square Point via Quadratic Obstructions}
\author{Edson Soares}
\date{Expanded preprint}
\subjclass[2020]{14D15, 14Q20, 68Q17, 13P10}
\keywords{Brent equations, fast matrix multiplication, local dimension,
quadratic obstructions, computational algebraic geometry}
\begin{document}
\begin{abstract}
We determine the local Krull dimension at the ordered recursive
Strassen-square point of the Brent scheme of 49-summand decompositions
of the $4\times4$ matrix multiplication tensor. The local dimension is
197 over $\overline{\F}_3$ and in characteristic zero, and the Zariski
tangent space has dimension 251. Two exact recursive deformation
families give the lower bound. For the upper bound, an exact
finite-field certificate produces 54 quadratic initial forms with
pairwise coprime leading monomials, selected by bipartite matching.
A saturated Jacobian over $\Z_{(3)}$ lifts the regular sequence to
characteristic zero. The computation also certifies two smooth formal
branches of dimension 197 with intersection dimension 143, and the
characteristic-three local Hilbert function begins $(1,251,30978)$.
\end{abstract}
\maketitle

\section{The point and the local claim}

For a field $k$, let $M_n\in(k^{n^2})^{\otimes3}$ have entries
$(M_n)_{(i,h),(h,j),(i,j)}=1$ for $0\le i,h,j<n$, with all other
entries zero. Ordered factor triples are written $(u_s,v_s,w_s)$,
where the coordinates of the first two factors specify input linear
forms and those of the third specify the output matrix. For $n=4$
and $r=49$ the Brent map has $2352=3\cdot49\cdot16$ variables and
$4096=16^3$ equations
\begin{equation}\label{eq:brent}
 \Phi_{abc}=\sum_{s=1}^{49}u_{sa}v_{sb}w_{sc}-(M_4)_{abc}=0.
\end{equation}
Write $\Br_k$ for the affine scheme defined by these equations and
$V_k(4,4,4\mid49)=(\Br_k)_{\mathrm{red}}$. Their local Krull dimensions
agree. We use the scheme local ring for the Hilbert function and
initially for the tangent computation.

Here is the particular integral point used throughout. The rows
of the three arrays are the seven Strassen summands in row-major
coordinates \cite{strassen}.
\[
 U=\begin{pmatrix}
 1&0&0&1\\0&0&1&1\\1&0&0&0\\0&0&0&1\\1&1&0&0\\-1&0&1&0\\0&1&0&-1
 \end{pmatrix}\;
 V=\begin{pmatrix}
 1&0&0&1\\1&0&0&0\\0&1&0&-1\\-1&0&1&0\\0&0&0&1\\1&1&0&0\\0&0&1&1
 \end{pmatrix}\;
 W=\begin{pmatrix}
 1&0&0&1\\0&0&1&-1\\0&1&0&1\\1&0&1&0\\-1&1&0&0\\0&0&0&1\\1&0&0&0
 \end{pmatrix}.
\]
For each of the three arrays, regard each row as a $2\times2$ matrix.
The $49$ rows of the corresponding array at $p$ are the row-major
vectorizations of the Kronecker products of row matrices, ordered
lexicographically by their pair of summand indices. Explicitly,
coordinate $((2i+i'),(2j+j'))$ of such a product is the product of
coordinates $(i,j)$ and $(i',j')$. Squaring the seven-summand tensor
identity and using this shuffle gives $p\in\Br(\Z)$. The verifier
rechecks the full identity over the integers.

\begin{theorem}\label{thm:main}
Let $k=\overline{\F}_3$ or let $k$ have characteristic zero. At the
recursive point $p$,
\[
 \rk D\Phi_p=2101,\qquad
 \dim T_p\Br_k=251,\qquad
 \dim_p V_k(4,4,4\mid49)=197.
\]
The continuous isotropy and term-rescaling orbit has dimension 143
at $p$. There are at least two distinct smooth 197-dimensional
formal branches through $p$, with scheme-theoretic formal
intersection of dimension 143. Over $\overline{\F}_3$ the local
Hilbert function of $\Br_k$ starts with $(1,251,30978)$.
\end{theorem}

The recursive square and the AlphaTensor catalogue give markedly
different constructions at the same length, $49=7^2$. A rational
48-term decomposition is already known \cite{dps}. Our question is
what the geometry at the classical recursive point says about how
49-term decompositions vary.

At a singular point, the Jacobian can leave directions that do not
integrate to deformations. For example, at the origin of $xy=0$ the
tangent space is two-dimensional, but the two coordinate axes are
the only branches and the local dimension is one. The recursive
construction makes this
failure concrete: two distinct ways of varying its two levels already
span the full tangent space. This raises a sharper question than the
rank of the Jacobian: how many parameters can coexist in an actual
family? The exact dimension and the intersection of the two families
answer part of that question.

Li, Bao and Zhang study local dimensions via Jacobian and orbit bounds
\cite{lbz}. Li, Zhang and Ke use deflation to treat singular Brent
solutions \cite{lzk}. Li, Wang and Hu report Jacobian nullity 197
and numerical evidence of smoothness for selected 49-summand
AlphaTensor points \cite[Section~5.2.2]{lwh}. Their selected example
has orbit tangent dimension 143, leaving 54 directions modulo the
orbit. Here 197 is the exact
local dimension at the recursive point, where the nullity is 251.
We compare the displayed families with the AlphaTensor catalogue in
Section~\ref{sec:alphatensor}.
Brysiewicz and Sinn \cite{bs} give general upper-bound tests from
quadratic obstructions. Our upper bound uses a certificate adapted to
the mixed structure of this particular point. Combining that bound
with the exact block motions yields equality.

\section{Two exact families and the tangent split}

Put $J=D\Phi_p$. Write $A_s,B_s,C_s$ for the factor matrices at a
summand, with $C_s$ in output order. The continuous group $G$ consists
of $\mathrm{GL}_4^3$ and two independent scalings per summand. For
$P,Q,R$ in the three general linear groups the factor maps are
\[
 A_s\longmapsto P^t A_s Q^{-t},\quad
 B_s\longmapsto Q^t B_s R^{-t},\quad
 C_s\longmapsto P^{-1} C_s R .
\]
These maps evaluate the original scheme on $(PXQ^{-1},QYR^{-1})$
and undo the corresponding output change. The term scalings
multiply the three factors by scalars with product~1.
Scalar $P,Q,R$ can be compensated by the term scalings, giving a
three-dimensional kernel in the $146$-dimensional parameter group.

The first recursive family applies a separate $\mathrm{GL}_2^3$
change of basis of the inner Strassen scheme for each fixed outer
summand. Each altered block still represents the tensor product of
that outer summand with $M_2$, with the same coordinate shuffle
used to define $p$. The second family exchanges the roles
of the tensor-product indices. Compose both constructions
separately with the same global group action. This gives two
morphisms from open subsets of smooth parameter spaces to the
Brent scheme, both passing through $p$.

We now derive the dimensions of these families by hand. Work in
characteristic different from two. Transpose the output factor, writing
$D_i=C_i^t$, so the three infinitesimal actions on the base scheme
have the cyclic form
\[
 \delta A_i = aA_i-A_i b,\qquad
 \delta B_i = bB_i-B_i c,\qquad
 \delta D_i = cD_i-D_i a.
\]
Here the three matrices $a,b,c$ are arbitrary $2\times2$ matrices.
Let $H=\mathfrak{sl}_2(k)$, the space of trace-zero matrices, so
$\operatorname{Mat}_2(k)=kI\oplus H$.

\begin{lemma}\label{lem:base}
The three copies of $H$ act injectively on the seven ordered triples
of factor lines of Strassen's decomposition, at the infinitesimal level.
\end{lemma}
\begin{proof}
Suppose each factor derivative is a multiple of that factor in every
summand. The first summand has all three factors $I$, so $a-b$,
$b-c$ and $c-a$ are scalar. Their traces are zero, giving $a=b=c=h$.
The first-mode factors include $E_{11}$ and $E_{11}+E_{12}$.
If $[h,E_{11}]$ is a multiple of $E_{11}$, the two off-diagonal
entries of $h$ vanish. Write $h=\operatorname{diag}(t,-t)$.
Then $[h,E_{11}+E_{12}]=2t E_{12}$ is a multiple of
$E_{11}+E_{12}$ only for $t=0$.
\end{proof}

Let $S$ be the term-rescaling tangent space, of dimension $2\cdot49=98$.
Let $L$ and $R'$ be the images of seven independent triples in $H^3$,
applied in the inner and outer positions, respectively. Consider any
one mode with factors $F_i\otimes F_j$ at grid position $(i,j)$. Modulo
that factor line, a derivative from $L$ lies in
$kF_i\otimes\operatorname{Mat}_2$ and one from $R'$ lies in
$\operatorname{Mat}_2\otimes kF_j$. Their only possible common
direction before taking the quotient is the factor line itself.
A relation among $S,L,R'$ therefore makes every individual block
motion trivial on the factor lines. Lemma~\ref{lem:base} kills each
block parameter. Thus $S+L+R'$ is a direct sum of dimensions
$98,63,63$.

For the global action use the decomposition
\[
 \operatorname{Mat}_4
   =(kI\otimes kI)\oplus (H\otimes kI)
     \oplus(kI\otimes H)\oplus(H\otimes H).
\]
The scalar part lands in $S$. The two single-factor parts give a
uniform block motion in each of $L$ and $R'$, of dimension nine each.
We claim that the three copies of $H\otimes H$ give 27 independent
directions modulo $S+L+R'$. To see this, project each factor derivative
at $(i,j)$ to
\[
 (\operatorname{Mat}_2/kF_i)\otimes
 (\operatorname{Mat}_2/kF_j).
\]
This projection annihilates $S+L+R'$. At the identity position
$(0,0)$, where all factor matrices are $I\otimes I$, vanishing in
all three modes forces the three mixed global generators to be
equal, since $H$ maps isomorphically to $\operatorname{Mat}_2/kI$.
Denote that common generator by $K\in H\otimes H$. At $(i,0)$ in
the first mode, the projected derivative is the projection of
$[K,A_i\otimes I]$. Write $K=\sum_\nu h_\nu\otimes e_\nu$ for
a basis $e_\nu$ of $H$. Its image vanishes only if each
$[h_\nu,A_i]$ is a multiple of $A_i$ for all $i$. The two
first-mode factors in Lemma~\ref{lem:base} then give $h_\nu=0$.
This proves the claim.

The global and branch tangent dimensions now follow:
\[
\begin{array}{c|c|c}
\text{Space} & \text{Dimension calculation} & \text{Dimension}\\ \hline
T_0 & 98+9+9+27 & 143\\
T_0+L,\ T_0+R' & 143+(63-9) & 197\\
T_0+L+R' & 98+63+63+27 & 251 .
\end{array}
\]
The two branch tangent spaces intersect precisely in $T_0$.
The global orbit has dimension 143, by the central kernel bound and
the differential rank just proved. The finite rank computation is
needed to show that these analytic tangent directions exhaust the
kernel: a rank-2101 minor modulo~3 gives the matching upper bound on
the tangent dimension in characteristic three and in characteristic zero.
Choosing complements $X$ and $Y$ from $L$ and $R'$ gives
\begin{equation}\label{eq:split}
 \ker J=T_0\oplus X\oplus Y,
 \qquad (\dim T_0,\dim X,\dim Y)=(143,54,54).
\end{equation}
The verifier uses integral generators in this order to select a particular
basis for the quadrics. The dimension argument is independent of that
basis selection.

Select 143 independent group parameters and then 54 independent
recursive parameters in each family. Their common restriction when
the recursive parameters are zero is the same group
parameterization. In each case 197 ambient coordinate functions
have invertible derivative with respect to these 197 parameters.
The formal inverse function theorem gives surjections from the
completed Brent local ring onto two power series rings in 197
variables. This proves the lower bound of 197 and gives embedded
smooth formal subgerms with distinct tangent images $T_0+X$ and
$T_0+Y$.

\section{The quadratic certificate in characteristic three}

In the expansion of $\Phi(p+v)$, let $Q_2(v)$ denote the degree-two
part in $v$. For a single summand with variation $(a,b,c)$ its
contribution is
\[
 a\otimes b\otimes w + a\otimes v_s\otimes c
  +u\otimes b\otimes c .
\]
For $\lambda J=0$ the restriction of $\lambda Q_2$ to $\ker J$
vanishes on $T_0+X$ and on $T_0+Y$. Indeed every direction in either
space lifts to a formal arc in the corresponding exact family, so
the order-two contribution lies in $\im J$. Thus the restriction
contains only the mixed monomials $x_i y_j$ in the two complements.

In the frozen bases, $M$ is the $4096\times2916$ matrix whose
$(i,j)$-column is
$Q_2(X_i+Y_j)-Q_2(X_i)-Q_2(Y_j)$.
The certificate gives 54 explicit rows $\lambda_i$ over $\F_3$ with
$\lambda_iJ=0$. The verifier reconstructs all tensor columns and
checks every coefficient of the 54 forms $q_i=\lambda_iM(x\otimes y)$.
In a common lexicographic order, with $x_0>\cdots>x_{53}$ followed
by $y_0>\cdots>y_{53}$, the monic leading terms are
\begin{equation}\label{eq:leading}
  \operatorname{LM}(q_i)=x_i y_{\pi(i)},\qquad 0\le i<54
\end{equation}
for a permutation $\pi$. In particular, these terms form a perfect
matching of the two variable sets. The rows can be found by
scanning the mixed monomial columns in order modulo $\im J$.
An independent column admits a functional annihilating the
earlier columns and $J$. Bipartite matching selects 54 such
positions. The supplied certificate is checked directly and
does not require trusting the selection algorithm.

By Buchberger's product criterion the $q_i$ form a Gr\"obner
basis of their homogeneous ideal. Its initial ideal is generated
by the 54 pairwise coprime monomials in \eqref{eq:leading}. Those
monomials generate a complete intersection of height 54. Hence
the homogeneous ideal has height 54 and its 54 generators form
a regular sequence in the polynomial ring, which is
Cohen--Macaulay. This holds also with the 143 common tangent
coordinates adjoined.

\begin{lemma}\label{lem:initial}
Let $A=k[[z_1,\ldots,z_N]]$ with maximal ideal $\mathfrak m$, and let
$I\subseteq A$. For a nonzero $f\in A$, write $f^*$ for its lowest
nonzero homogeneous part. If $I$ contains $f_1,\ldots,f_s$ with
$f_1^*,\ldots,f_s^*$ a regular sequence, then $\dim A/I\le N-s$.
For $I=(f_1,\ldots,f_s)$ the associated graded ring is
$k[z_1,\ldots,z_N]/(f_1^*,\ldots,f_s^*)$.
\end{lemma}
\begin{proof}
For any $I$, the graded map $\operatorname{gr}_{\mathfrak m}A\to
\operatorname{gr}_{\mathfrak m}(A/I)$ has kernel generated by the
initial forms of elements of $I$. Indeed, a class of degree $d$
in its kernel has a representative in
$\mathfrak m^d\cap(I+\mathfrak m^{d+1})$ and hence agrees with
an element of $I\cap\mathfrak m^d$ modulo $\mathfrak m^{d+1}$.
Taking the quotient by the $s$ specified initial forms therefore
gives a graded surjection onto $\operatorname{gr}_{\mathfrak m}(A/I)$.
The source has dimension $N-s$. The Hilbert--Samuel dimension
theorem gives $\dim\operatorname{gr}_{\mathfrak m}(A/I)=\dim A/I$.

For the equality of graded rings in the stated special case, proceed
one series at a time. In the quotient by the preceding series,
the next $f_i^*$ is a nonzerodivisor in the graded ring by induction.
For every nonzero element $g$ of that quotient, the product of the
two initial forms is nonzero, so $(gf_i)^*=g^*f_i^*$. Thus the initial
ideal of $(f_i)$ in this quotient is precisely $(f_i^*)$. Induction
gives the asserted presentation.
\end{proof}

Split the ambient coordinate space into $\ker J$, with coordinates
$z=(c,x,y)$, and a normal complement, with coordinates $n$.
Choose 2101 equations $F$ whose derivative in $n$ is an invertible
matrix $D$. Then $F(z,n)=Dn+O((z,n)^2)$. The formal map
$(z,n)\mapsto(z,F(z,n))$ has linear part $\operatorname{diag}(1,D)$
and has a formal inverse: solving for each successive homogeneous
part of that inverse requires only multiplication by $D^{-1}$
in the normal coordinates. Setting $F=0$ gives the unique normal
series $n(z)$ of order at least two. The completed Brent local
ring is $k[[z]]/I$, where $I$ is generated by the original
equations after substituting $n(z)$.

In particular $I$ contains $f_i(z)=\lambda_i\Phi(p+z+n(z))$.
The linear contribution dies under $\lambda_i$ because $\lambda_i J=0$.
Since $n(z)$ starts in degree two, it contributes only in degree
three and higher to the quadratic and cubic parts of the expansion.
Thus $f_i^*=q_i$. The dimension inequality in
Lemma~\ref{lem:initial} gives $\dim k[[z]]/I\le251-54=197$.
Together with the lower bound this proves the characteristic-three
dimension claim.

\section{Transfer to characteristic zero}

Let $R_0=\Z_{(3)}$. The Jacobian has a rank-2101 minor that is a
unit over $R_0$. The 251 independent integral tangent generators
in \eqref{eq:split} remain independent over $\Q$ and annihilate
$J$ integrally, so the rank over $\Q$ is also 2101. Invertible
row and column operations over $R_0$ therefore split $J$ as an
identity block and a zero block: the remaining Schur complement
vanishes over $\Q$, hence over $R_0$. Both Jacobian kernels are
free direct summands and reduce onto the corresponding kernels
modulo~3. The integral tangent generators selected in
\eqref{eq:split} give a basis over $R_0$ by reduction.

Lift the 54 rows $\lambda_i$ into the left kernel over $R_0$.
The normal elimination uses the same unit minor and works
formally over $R_0$. The quadratic parts $\widetilde q_i$ of the
lifted residual equations reduce to $q_i$. In the local polynomial
ring $R_0[z]_{(3,z)}$, the sequence $3,\widetilde q_1,\ldots,
\widetilde q_{54}$ is regular by the regularity of its reduction
after the first element. Permute the sequence to put $3$ last
\cite[Section~10.68]{stacks}. Localizing at $(z)$ gives
$\Q[z]_{(z)}$, where the 54 quadrics remain a regular sequence
and generate a proper ideal. Flat completion gives the same
property in $\Q[[z]]$. Thus Lemma~\ref{lem:initial} gives the
same upper bound in characteristic zero. The integral family
minors give the lower bound in characteristic zero, including
over $\mathbb C$.

\section{Formal branches and Hilbert function}

Each smooth formal subgerm constructed above gives a prime of
dimension 197 in the completed local ring. These primes are
minimal by the dimension bound. Their distinct tangent images
make the primes distinct. The intersection contains the shared
143-dimensional smooth orbit subgerm and its tangent space is
the intersection of the displayed tangent images, of dimension
143. This proves that the intersection itself is regular of dimension
143. In the 251-variable formal ambient space one can straighten the
two subgerms simultaneously. Indeed, if their parameterizations are
$f_L(c,x)$ and $f_R(c,y)$ with common restriction $f_0(c)$, the sum
$f_L(c,x)+f_R(c,y)-f_0(c)$ has invertible derivative. In the resulting
coordinates the two subgerms are $x=0$ and $y=0$.
The two smooth subgerms factor through the reduction, and their
tangents span a 251-dimensional space. Thus the tangent space
of the reduced variety also has dimension 251 and the reduced
variety is singular at $p$.

In characteristic three the reconstructed matrices give
$\rk[J\mid M]=2749$. After eliminating the normal coordinates,
the degree-two initial part of the full residual ideal is the
image of the mixed forms from the entire left kernel of $J$.
It has dimension $2749-2101=648$. Products by positive-degree
multipliers do not contribute in degree two, so the local
Hilbert value there is $\binom{252}{2}-648=30978$.
The reduced union of the two displayed subgerms would impose all
$54^2=2916$ mixed monomials, with degree-two Hilbert value $28710$.
The measured value $30978$ thus gives a dichotomy in characteristic
three: the completed local ring is nonreduced, or there is an
additional minimal prime. The full Hilbert series and branch count
remain open.

\section{Relation to the AlphaTensor schemes}\label{sec:alphatensor}

The common dimension 197 has a concrete explanation on the recursive
families. The closure of either family image has dimension 197:
its differential has rank 197 at $p$, and its dimension is bounded
above by the local dimension of the Brent scheme at $p$. Each closure
is an irreducible component of the reduced variety through $p$.
This argument allows the two global closures to coincide, as can
happen for branches at a singular point.

We exhibit smooth points of the full Brent scheme on both displayed
families. In each orientation the accompanying verifier constructs
a point using integer unimodular changes of basis, checks all 4096
integer equations and finds Jacobian rank 2155 modulo three.
Because each point lies on a 197-dimensional component, its tangent
dimension is at least 197. The calculation gives the reverse inequality,
both in characteristic three and in characteristic zero. Equality of
embedding and local dimension proves smoothness at these points.
This shows how a singular 49-term scheme with tangent dimension 251
can be connected by explicit deformations to smooth schemes with tangent
dimension 197, the same dimension reported for many AlphaTensor
examples \cite[Section~5.2.2]{lwh}. Identifying which catalogue points
are on which component would require additional information.

We can nevertheless rule out membership for some catalogue points.
Throughout either displayed family, the multiset of triples of ordinary
matrix ranks is
\[
 36(1,1,1),\qquad 12(2,2,2),\qquad 1(4,4,4).
\]
The base algorithm has six triples of rank one and a single triple
of rank two. Tensor products multiply these ranks, and the permitted
changes are left and right multiplication by invertible matrices.
In the closures, ranks can only decrease at each labeled summand.
Consequently, even allowing summand and mode permutations, at most 13
summands can have any mode rank above one, and at most one summand
can have any mode rank above two.

A second necessary condition applies to catalogue points with the
same rank profile. At each of the six rank-one outer summands, the
group of seven terms contains a unique rank-two term, obtained by
taking the inner identity summand. In each mode, its row and column
spaces contain the row and column spaces of all six other terms in
that group. This follows from the inclusion of the column and row
spaces in a fixed one-dimensional space tensored with $k^2$.
Left and right changes preserve the inclusions. They persist at
closure points with the same profile because the relevant minors of
the stacked matrices vanish. Indeed, equality of the total rank sum
in the two profiles forces equality at every label, since each rank
is individually bounded above by its original value. In particular
each anchor stays of rank two.
The other recursive orientation gives the same condition with the
index roles exchanged. The filter checks that at least six rank-two
terms can each accommodate six rank-one terms, and that each rank-one
term has at least one such candidate anchor. This test is necessary
for an actual partition.

We applied these tests to the 14,236 public integral factorizations
of Fawzi et al.\ \cite{alpha}, using exact row reduction. In both
characteristics the rank bounds exclude 87 points from the closures
of the two displayed families, and the row- and column-space test
excludes two further points with the same profile. The remaining
14,147 pass these necessary tests. The supplementary script records
the exact input revision, checks its SHA-256 hash, and uses a prime
larger than the Hadamard bound for every relevant minor when testing
characteristic zero. These exclusions concern the closures of the
displayed families. A catalogue point outside these closures could
still be on another component through $p$.

The open parameterizations preserve every nonzero factor and retain
all 49 summands. Thus they give deformations of a fixed length.
A shorter decomposition would require a suitably degenerate limit
or another construction. The existence of a length-reducing limit
remains open here.

\section{Reproducibility}

The \href{https://github.com/edsonphx/brent-strassen-square}{public
repository} includes the signed factor arrays, integral tangent
bases, sparse functionals, quadratic forms and SHA-256 manifest.
Run \texttt{bash verify.sh} with Python and NumPy to check the tensor
identity, tangent equations, Jacobian and augmented ranks, and
the 54 functionals and matching positions. The verifier derives
the mixed columns by signed polarization of integer contractions.
Pass \texttt{--full} to check all 28,710 non-mixed monomials per
certified form and verify 40 exact family points, including the
two smooth witnesses. A separate script runs the catalogue census.

The original five-page version and its verifier are archived as
\href{https://github.com/edsonphx/brent-strassen-square/releases/tag/v1.0.0}{GitHub release \texttt{v1.0.0}}
and at \href{https://doi.org/10.5281/zenodo.24449376}{DOI 10.5281/zenodo.24449376}.

\begin{thebibliography}{9}
\bibitem{lbz} X. Li, Y. Bao and L. Zhang,
On the local dimensions of solutions of Brent equations,
\emph{Linear Algebra Appl.} 708 (2025), 489--512.
\href{https://doi.org/10.1016/j.laa.2024.12.011}{doi:10.1016/j.laa.2024.12.011}.
\bibitem{lzk} X. Li, L. Zhang and Y. Ke,
Deflation conjecture and local dimensions of Brent equations,
\emph{Exp. Math.} 34 (2025), 488--501.
\href{https://arxiv.org/abs/2310.11686}{arXiv:2310.11686}.
\bibitem{lwh} X. Li, Y. Wang and S. Hu,
Substitution and quotient of the isotropy group action,
\href{https://arxiv.org/abs/2607.15069v2}{arXiv:2607.15069v2} (2026).
\bibitem{dps} J.-G. Dumas, C. Pernet and A. Sedoglavic,
A non-commutative algorithm for multiplying $4\times4$ matrices using
48 non-complex multiplications,
\href{https://arxiv.org/abs/2506.13242}{arXiv:2506.13242} (2025).
\bibitem{bs} T. Brysiewicz and R. Sinn,
\href{https://arxiv.org/abs/2609.02818}{Local dimension testing via
quadratic obstructions} (2026).
\bibitem{alpha} A. Fawzi et al.,
\href{https://doi.org/10.1038/s41586-022-05172-4}{Discovering faster
matrix multiplication algorithms with reinforcement learning},
\emph{Nature} 610 (2022), 47--53.
\bibitem{stacks} The Stacks Project Authors,
\emph{The Stacks Project}, Section~10.68,
\href{https://stacks.math.columbia.edu/tag/0AUH}{Tag 0AUH}.
\bibitem{strassen} V. Strassen,
Gaussian elimination is not optimal,
\emph{Numer. Math.} 13 (1969), 354--356.
\href{https://doi.org/10.1007/BF02165411}{doi:10.1007/BF02165411}.
\end{thebibliography}
\end{document}
